% Vietnamese translation for OI Public Library
% Source-PDF: IOI中国国家候选队论文/2009/梅诗珂/信息学竞赛中概率问题求解初探.pdf
\documentclass[11pt]{article}
\usepackage{oipl-vn}

\newcommand{\Prob}{\operatorname{P}}
\newcommand{\E}{\operatorname{E}}
\newcommand{\dd}{\,\mathrm{d}}
\newcommand{\R}{\mathbb{R}}
\newcommand{\Vol}{\operatorname{V}}
\newcommand{\Total}{\mathit{Total}}

\title{Bước vào thế giới xác suất: Bàn sơ về cách giải bài toán xác suất trong thi Tin học}
\author{Mei Shike\\Trường Trung học số 1 Hợp Phì, An Huy}
\date{2009}

\begin{document}
\maketitle

\origtitle{Bước vào thế giới xác suất: Bàn sơ về cách giải bài toán xác suất trong thi Tin học}
\sourcepath{IOI China candidate team papers / 2009 / Mei Shike}

\begin{abstract}
Việc thiết kế nhiều thuật toán trong Tin học có liên quan tới xác suất. Trong
thi Tin học, các bài toán yêu cầu tính xác suất hoặc kỳ vọng cũng chiếm một
tỉ trọng đáng kể và thường khá khó. Bài viết dùng các tính chất của tổ hợp,
phân tích sai số, chuyển sang phần bù, chia hàm theo đoạn và một số kỹ thuật
khác để giải bốn ví dụ, từ đó tổng kết các đặc điểm chung của bài toán xác
suất và các chiến lược tương ứng.
\end{abstract}

\paragraph{Từ khóa.}
Xác suất; biến ngẫu nhiên; liên tục; rời rạc; mật độ xác suất; tích phân.

\section{Kiến thức cơ sở}

\subsection{Không gian mẫu, sự kiện và xác suất}

\term{Không gian mẫu} \(S\) là một tập hợp; các phần tử của nó được gọi là
\term{sự kiện sơ cấp}. Một tập con của không gian mẫu được gọi là
\term{sự kiện}. Theo định nghĩa, mọi sự kiện sơ cấp loại trừ lẫn nhau.

Nếu có một ánh xạ từ các sự kiện vào tập số thực, ký hiệu là \(\Prob\{\cdot\}\),
thỏa mãn:
\begin{enumerate}
  \item Với mọi sự kiện \(A\), \(\Prob\{A\}\ge 0\).
  \item \(\Prob\{S\}=1\).
  \item Với hai sự kiện loại trừ lẫn nhau, \(\Prob\{A\cup B\}=\Prob\{A\}+\Prob\{B\}\).
\end{enumerate}
thì ta gọi \(\Prob\{A\}\) là xác suất của sự kiện \(A\). Ba điều trên gọi là
các tiên đề xác suất.

\subsection{Biến ngẫu nhiên}

Nếu với mỗi sự kiện sơ cấp \(e\) trong không gian mẫu \(S\) có đúng một số thực
\(X(e)\) tương ứng, thì \(X=X(e)\) được gọi là một biến ngẫu nhiên trên
\(S\). Hai loại thường gặp là biến ngẫu nhiên rời rạc và biến ngẫu nhiên liên
tục.

\subsubsection{Biến ngẫu nhiên rời rạc và phân bố xác suất}

Biến ngẫu nhiên có miền giá trị là hữu hạn hoặc vô hạn đếm được các số thực
được gọi là biến ngẫu nhiên rời rạc. Giả sử xác suất để biến ngẫu nhiên rời
rạc \(X\) nhận giá trị \(x_k\) là \(p_k\), với \(k=1,2,\ldots\). Khi đó toàn bộ
các giá trị của \(X\) cùng xác suất tương ứng được gọi là phân bố xác suất của
\(X\), ký hiệu
\[
  \Prob\{X=x_k\}=p_k,\qquad k=1,2,\ldots.
\]
Các phân bố rời rạc thường gặp gồm phân bố hai điểm, nhị thức, hình học, siêu
bội và Poisson.

\subsubsection{Biến ngẫu nhiên liên tục và phân bố xác suất}

Nếu \(X\) là biến ngẫu nhiên nhận giá trị liên tục trên trục số thực hoặc trên
một khoảng, đặt hàm phân bố xác suất của \(X\) là
\[
  F(x)=\Prob\{X\le x\}.
\]
Nếu tồn tại một hàm không âm, khả tích \(f(x)\), sao cho với mọi \(x\),
\[
  F(x)=\int_{-\infty}^{x} f(t)\dd t,
\]
thì \(X\) được gọi là biến ngẫu nhiên liên tục, và \(f(x)\) được gọi là hàm
mật độ xác suất của \(X\). Cần chú ý rằng mật độ xác suất không phải là xác
suất. Các phân bố liên tục thường gặp gồm phân bố đều, chuẩn và mũ.

\subsubsection{Vector ngẫu nhiên liên tục và phân bố của nó}

Nếu \(X_1,X_2,\ldots,X_N\) đều là biến ngẫu nhiên liên tục, thì
\((X_1,X_2,\ldots,X_N)\) được gọi là vector ngẫu nhiên \(N\) chiều. Hàm phân bố
xác suất của nó là
\[
  F(x_1,x_2,\ldots,x_N)
  =\Prob\{X_1\le x_1,X_2\le x_2,\ldots,X_N\le x_N\}.
\]
Nếu tồn tại một hàm không âm, khả tích \(f(x_1,x_2,\ldots,x_N)\), sao cho
\[
  F(x_1,x_2,\ldots,x_N)
  =\int_D f(x_1,x_2,\ldots,x_N)\dd x_1\dd x_2\cdots\dd x_N,
\]
trong đó vế phải là tích phân \(N\) lớp trên miền thích hợp \(D\), thì
\(f(x_1,x_2,\ldots,x_N)\) được gọi là hàm mật độ xác suất đồng thời của vector
ngẫu nhiên \((X_1,X_2,\ldots,X_N)\). Nếu \(X_1,X_2,\ldots,X_N\) độc lập với
nhau và lần lượt có các mật độ \(f_1(x_1),f_2(x_2),\ldots,f_N(x_N)\), thì
\[
  f(x_1,x_2,\ldots,x_N)
  =f_1(x_1)f_2(x_2)\cdots f_N(x_N)
\]
là một mật độ đồng thời hợp lệ.

\subsection{Kỳ vọng toán học}

\subsubsection{Kỳ vọng của biến ngẫu nhiên rời rạc}

Giả sử biến ngẫu nhiên rời rạc \(X\) có luật phân bố
\(\Prob\{X=x_k\}=p_k\), với \(k=1,2,\ldots\). Nếu
\[
  \sum_{k=1}^{\infty}|x_kp_k|
\]
tồn tại, thì
\[
  \sum_{k=1}^{\infty}x_kp_k
\]
được gọi là kỳ vọng toán học của \(X\), thường gọi tắt là kỳ vọng, ký hiệu
\(\E(X)\).

\subsubsection{Kỳ vọng của biến ngẫu nhiên liên tục}

Giả sử biến ngẫu nhiên liên tục \(X\) có hàm mật độ xác suất \(f(x)\). Nếu
tích phân suy rộng
\[
  \int_{-\infty}^{+\infty}|xf(x)|\dd x
\]
hội tụ, thì
\[
  \int_{-\infty}^{+\infty}xf(x)\dd x
\]
được gọi là kỳ vọng của biến ngẫu nhiên liên tục \(X\), ký hiệu \(\E(X)\).

\subsection{Tích phân}

Giả sử hàm \(f\) được định nghĩa trên đoạn đóng \([a,b]\). Gọi một phân hoạch
\(\pi\) của \([a,b]\) là
\[
  (x_0=a,x_1,x_2,\ldots,x_n=b),\qquad
  x_0<x_1<x_2<\cdots<x_n,
\]
và đặt
\[
  \|\pi\|=\max_{1\le i\le n}(x_i-x_{i-1}).
\]
Lấy tùy ý \(\xi_i\in[x_{i-1},x_i]\) làm đại diện cho đoạn con
\([x_{i-1},x_i]\). Nếu tồn tại một số \(A\) sao cho với mọi
\(\varepsilon>0\) đủ nhỏ luôn tồn tại \(\delta>0\), và miễn là
\(\|\pi\|<\delta\) thì
\[
  \left|\sum_{i=1}^{n} f(\xi_i)(x_i-x_{i-1})-A\right|<\varepsilon
\]
không phụ thuộc vào cách chọn \(\xi_i\), thì \(f\) được gọi là khả tích trên
\([a,b]\), và \(A\) được gọi là tích phân xác định của \(f\) trên \([a,b]\),
ký hiệu
\[
  \int_a^b f(x)\dd x.
\]
Một kết luận thực dụng là mọi hàm liên tục trên một đoạn đóng đều khả tích.

Công thức Newton--Leibniz: nếu hàm liên tục \(f\) được định nghĩa trên
\([a,b]\), thì
\[
  \int_a^b f(x)\dd x=F(b)-F(a),
\]
trong đó \(F(x)\) là một nguyên hàm bất kỳ của \(f(x)\).

Dưới đây ta lần lượt phân tích hai ví dụ về biến ngẫu nhiên rời rạc và hai ví
dụ về biến ngẫu nhiên liên tục.

\section{Bài toán về biến ngẫu nhiên rời rạc}

Bài toán về biến ngẫu nhiên rời rạc đã nhiều lần xuất hiện trong thi Tin học,
chẳng hạn \emph{Congcong và Coco} của NOI 2005, \emph{Chiếc túi thần kỳ} của
NOI 2006, hay \emph{Sắp lịch thi đấu} của NOI 2008. Loại bài này thường không
đòi hỏi thuật toán quá phức tạp; chỉ cần chứng minh được định lý then chốt thì
bài toán đã được giải quyết. Muốn xử lý loại bài này, cần nắm khá vững các
tính chất của những phân bố rời rạc thường gặp.

Ta bắt đầu bằng một ví dụ về xác suất thắng tối đa.

\subsection{Ví dụ 1: LastMarble}

\paragraph{Tóm tắt đề bài.}
Trong một túi có \(\mathit{red}\) viên bi đỏ và \(\mathit{blue}\) viên bi
xanh. Hai người chơi luân phiên lấy bi từ túi; mỗi lượt một người có thể lấy
\(1\), \(2\) hoặc \(3\) viên, nhưng trước khi lấy ra thì không biết màu của
chúng. Sau khi bi được lấy ra, màu của chúng được công bố cho cả hai người.
Người lấy viên bi đỏ cuối cùng sẽ thua. Biết rằng trước khi ván chơi bắt đầu,
có người đã lấy ngẫu nhiên \(\mathit{removed}\) viên bi và không ai biết màu
của chúng. Nếu hai người đều chơi tối ưu, xác suất thắng của người đi trước là
bao nhiêu?

Ràng buộc:
\[
  1\le \mathit{red},\mathit{blue}\le 100,\qquad
  0\le \mathit{removed}\le \mathit{red}-1.
\]

\paragraph{Phân tích.}
Khi \(\mathit{removed}=0\), đây là một bài quy hoạch động rất bình thường.
Đặt \(F(r,b)\) là xác suất thắng tối đa của người đang đối mặt với trạng thái
còn \(r\) viên đỏ và \(b\) viên xanh. Khi đó
\[
  F(0,b)=1\qquad (b\ge 0),
\]
và
\[
  F(r,b)=
  \max_{1\le m\le \min(r+b,3)}
  \left(
    \sum_{\substack{0\le p\le r,\ 0\le q\le b\\p+q=m}}
    \frac{\binom{r}{p}\binom{b}{q}}{\binom{r+b}{m}}
    \bigl(1-F(r-p,b-q)\bigr)
  \right).
  \tag{1}
\]
Ở đây
\[
  \frac{\binom{r}{p}\binom{b}{q}}{\binom{r+b}{m}}
\]
là xác suất lấy được \(p\) viên đỏ và \(q\) viên xanh. Khi đó
\(F(\mathit{red},\mathit{blue})\) là đáp án. Để xử lý
\(\mathit{removed}>0\), ta cần định lý sau.

\paragraph{Định lý 1.}
Trong một túi có \(\mathit{red}\) viên đỏ và \(\mathit{blue}\) viên xanh, nếu
lấy trước \(a\) viên rồi lấy tiếp \(b\) viên, thì phân bố xác suất của số bi
từng màu còn lại giống hệt trường hợp lấy trước \(b\) viên rồi lấy tiếp
\(a\) viên.

\paragraph{Chứng minh.}
Định lý này nhìn qua khá hiển nhiên, vì lấy \(a\) viên rồi lấy \(b\) viên phải
tương đương với lấy trực tiếp \(a+b\) viên. Ta chứng minh bằng đại số. Trước
hết,
\[
  \binom{n}{a}\binom{n-a}{b}
  =\binom{n}{a+b}\binom{a+b}{a}
  \qquad(a+b\le n),
\]
vì
\[
\begin{aligned}
  \binom{n}{a}\binom{n-a}{b}
  &=
  \frac{n!}{a!(n-a)!}\cdot
  \frac{(n-a)!}{b!(n-a-b)!} \\
  &=
  \frac{n!}{(a+b)!(n-a-b)!}\cdot
  \frac{(a+b)!}{a!b!}
  =\binom{n}{a+b}\binom{a+b}{a}.
\end{aligned}
\]
Giả sử sau khi lấy \(a\) viên rồi lấy \(b\) viên, tổng số bi đỏ đã bị lấy là
\(r\). Gọi \(a_1\) là số bi đỏ lấy được trong lần đầu. Vì các lượt lấy là độc
lập theo điều kiện số bi còn lại, xác suất của trường hợp này là
\[
\begin{aligned}
&\sum_{a_1=0}^{\min(a,r)}
  \frac{\binom{\mathit{red}}{a_1}\binom{\mathit{blue}}{a-a_1}}
       {\binom{\mathit{red}+\mathit{blue}}{a}}
  \cdot
  \frac{\binom{\mathit{red}-a_1}{r-a_1}
        \binom{\mathit{blue}-(a-a_1)}{b-(r-a_1)}}
       {\binom{\mathit{red}+\mathit{blue}-a}{b}}\\
&\quad =
  \sum_{a_1=0}^{\min(a,r)}
  \left(
  \frac{\binom{\mathit{red}}{r}\binom{\mathit{blue}}{a+b-r}}
       {\binom{\mathit{red}+\mathit{blue}}{a+b}}
  \cdot
  \frac{\binom{r}{a_1}\binom{a+b-r}{a-a_1}}
       {\binom{a+b}{a}}
  \right)\\
&\quad =
  \frac{\binom{\mathit{red}}{r}\binom{\mathit{blue}}{a+b-r}}
       {\binom{\mathit{red}+\mathit{blue}}{a+b}}.
\end{aligned}
\]
Ở dòng thứ hai, do một trong hai điều kiện \(a_1\le a\) và \(a_1\le r\) luôn
đúng theo miền lấy tổng, đẳng thức giữ nguyên. Kết quả cuối cùng nói rằng lấy
\(a\) viên rồi lấy \(b\) viên tạo ra đúng cùng phân bố với lấy trực tiếp
\(a+b\) viên. Vì vậy thứ tự lấy bi không ảnh hưởng tới kết quả cuối cùng.
\(\text{QED}\)

Đặt \(F(r,b)\) là xác suất thắng tối đa của người chơi hiện tại khi hiện còn
\(r\) viên đỏ, \(b\) viên xanh, và có \(\mathit{removed}\) viên đã bị lấy
trước nhưng không biết màu. Khi người chơi lấy \(m\) viên, theo định lý 1,
trạng thái mới hoàn toàn tương đương với việc người chơi lấy trước \(m\) viên
rồi để \(\mathit{removed}\) viên bị lấy đi. Tuy nhiên có một biên cần xử lý:
không thể bảo đảm \(\mathit{removed}\le r\), nên trong một vài trường hợp sau
khi người chơi lấy \(m\) viên thì người đó đã thua, không còn ván chơi phía
sau. Để xử lý trường hợp đặc biệt này, ta sửa nhẹ định nghĩa của \(F\):
\(F(r,b)\) là xác suất thắng tối đa của người chơi hiện tại, với điều kiện sau
khi \(\mathit{removed}\) viên bị lấy đi vẫn còn ít nhất một viên đỏ.

Ký hiệu \(\operatorname{Pro}(r,b,k)\) là xác suất rằng từ trạng thái có \(r\)
viên đỏ và \(b\) viên xanh, sau khi lấy \(k\) viên vẫn còn số bi đỏ lớn hơn
\(0\). Khi đó
\[
  \operatorname{Pro}(r,b,k)
  =
  \sum_{\substack{a<r\\0\le k-a\le b}}
  \frac{\binom{r}{a}\binom{b}{k-a}}{\binom{r+b}{k}}.
\]
Ta tính giá trị này bằng truy hồi:
\[
  \operatorname{Pro}(0,b,0)=0,\qquad
  \operatorname{Pro}(r,b,0)=1\quad(r>0),
\]
và
\[
  \operatorname{Pro}(r,b,k)
  =
  \frac{r}{r+b}\operatorname{Pro}(r-1,b,k-1)
  +\frac{b}{r+b}\operatorname{Pro}(r,b-1,k-1).
\]

Quay lại công thức (1). Trong các trường hợp mà từ \(r\) viên đỏ và \(b\) viên
xanh, lấy đi \(\mathit{removed}\) viên vẫn còn ít nhất một viên đỏ, xác suất
có điều kiện để sau khi lấy thêm \(p\) viên đỏ và \(q\) viên xanh rồi lấy
\(\mathit{removed}\) viên vẫn còn ít nhất một viên đỏ là
\[
  \frac{\operatorname{Pro}(r-p,b-q,\mathit{removed})}
       {\operatorname{Pro}(r,b,\mathit{removed})}.
\]
Do đó phương trình cuối cùng là
\[
  F(r,b)=0\qquad (r+b=\mathit{removed}),
\]
và với \(r+b>\mathit{removed}\),
\[
\begin{aligned}
F(r,b)=
\max_{1\le m\le \min(r+b,3)}
\sum_{\substack{p+q=m\\0\le p\le r,\ 0\le q\le b}}
&
\frac{\operatorname{Pro}(r-p,b-q,\mathit{removed})}
     {\operatorname{Pro}(r,b,\mathit{removed})}
\cdot
\frac{\binom{r}{p}\binom{b}{q}}{\binom{r+b}{m}}\\
&\cdot \bigl(1-F(r-p,b-q)\bigr).
\end{aligned}
\tag{2}
\]
Giá trị cần tìm là \(F(\mathit{red},\mathit{blue})\).

Quan sát công thức (2), ta thấy khác biệt duy nhất so với công thức (1) là
thừa số
\[
  \frac{\operatorname{Pro}(r-p,b-q,\mathit{removed})}
       {\operatorname{Pro}(r,b,\mathit{removed})}.
\]
Hai công thức giống nhau một cách đáng ngạc nhiên, nhưng cũng rất hợp lý. Nếu
không chứng minh được định lý trên thì lời giải có thể phức tạp hơn rất
nhiều. Điều này cho thấy vai trò quan trọng của chứng minh toán học trong bài
toán xác suất. Ở ví dụ tiếp theo, ta sẽ thấy tầm quan trọng của việc quan sát.

\subsection{Ví dụ 2: Randomness}

\paragraph{Tóm tắt đề bài.}
Có một bộ sinh số ngẫu nhiên, mỗi lần gọi trả về đều một số nguyên dương từ
\(1\) tới \(R\). Có \(N\) sự kiện; xác suất yêu cầu của sự kiện thứ \(i\) là
\[
  \frac{a_i}{b_i}.
\]
Hãy dùng bộ sinh số này để thiết kế một thiết bị kích hoạt sự kiện sao cho số
lần gọi kỳ vọng \(\mathit{Exp}\) là nhỏ nhất, rồi tính \(\mathit{Exp}\) với độ
chính xác \(10^{-7}\).

Ràng buộc:
\[
\begin{gathered}
  2\le R\le 1000,\qquad 1\le N\le 1000,\\
  1\le a_i\le b_i\le 1000\quad(1\le i\le N),\qquad
  \sum_{i=1}^{N}\frac{a_i}{b_i}=1.
\end{gathered}
\]
Tất cả các số đều là số nguyên dương.

\paragraph{Phân tích.}
Để quan sát quy luật, trước hết xét một ví dụ:
\[
  R=100,\qquad N=3,\qquad
  \frac{a_1}{b_1}=\frac{a_2}{b_2}=\frac{a_3}{b_3}=\frac13.
\]
Một chiến lược tối ưu là: nếu số sinh ra \(x\le 99\), thực hiện sự kiện
\(x\bmod 3\); nếu không, gọi lại bộ sinh số. Lần gọi sau xử lý theo cùng quy
tắc. Khi đó xác suất của mỗi sự kiện là
\[
  \sum_{i=1}^{\infty}\frac{33}{100^i}
  =\frac{33/100}{1-1/100}
  =\frac13,
\]
và
\[
  \mathit{Exp}=
  \sum_{i=0}^{\infty}\frac{1}{100^i}
  =\frac{1}{1-1/100}
  =\frac{100}{99}.
\]

Ta thấy ở lần gọi đầu tiên, gọi là bộ sinh tầng thứ nhất, có một số kết quả
được ánh xạ trực tiếp tới sự kiện; các kết quả còn lại mỗi cái tương ứng với
một lần gọi bộ sinh mới, gọi là bộ sinh tầng thứ hai. Quy nạp, bộ sinh tầng
\(i+1\) là lần gọi mới do một kết quả của bộ sinh tầng \(i\) gây ra.

Với sự kiện \(i\), giả sử ở tầng \(k\) có \(D(i,k)\) giá trị trả về làm sự
kiện \(i\) xảy ra. Khi đó nhất định có
\[
  \frac{a_i}{b_i}=\sum_{k=1}^{\infty}\frac{D(i,k)}{R^k}.
\]
Ta có thể giả sử \(D(i,k)<R\) với mọi \(k\). Nếu một \(D(i,k)\ge R\), ta giảm
\(D(i,k)\) đi \(R\) và tăng \(D(i,k-1)\) lên \(1\); đẳng thức vẫn đúng nhưng
kỳ vọng giảm. Vì vậy đây chính là biểu diễn cơ số \(R\) của
\(a_i/b_i\), nên các giá trị \(D(i,k)\) được xác định duy nhất.

Để tính \(\mathit{Exp}\), đặt \(H(i+1)\) là số kết quả ở tầng \(i\) dẫn tới
việc phải gọi thêm bộ sinh, và đặt \(H(1)=1\). Khi đó \(H(i)/R^i\) là đóng góp
của tầng \(i\) vào kỳ vọng, nên
\[
  \mathit{Exp}=\sum_{i=1}^{\infty}\frac{H(i)}{R^i}.
\]
Để tính \(H(i)\), chú ý rằng số kết quả ở tầng \(i\) dẫn tới gọi mới bằng
tổng số kết quả ở tầng đó trừ đi số kết quả dẫn tới sự kiện:
\[
  H(k+1)=H(k)R-\sum_{i=1}^{N}D(i,k).
\]

Ta không thể và cũng không cần tính mọi \(H(i)\). Với mọi \(m\ge 2\),
\[
  \sum_{k=m}^{\infty}\frac{H(k)}{R^k}
  <
  \sum_{k=m}^{\infty}\frac{NR}{R^k}
  =
  \frac{N}{R^{m-2}(R-1)}
  \le
  \frac{N}{R^{m-2}}.
\]
Bất đẳng thức này cho biết chỉ cần \(m\) đủ lớn, sai số của \(\mathit{Exp}\)
có thể nhỏ tùy ý, và tốc độ giảm sai số rất nhanh. Với độ chính xác mà đề bài
yêu cầu, tính tới tầng thứ \(50\) là đủ.

Nhìn lại quá trình phân tích, bước mấu chốt là quan sát ví dụ ban đầu. Từ ví
dụ đó ta mới có khái niệm tầng của bộ sinh ngẫu nhiên, và phần còn lại trở nên
tự nhiên. Điều này gợi ý rằng khi giải bài toán xác suất, xuất phát từ trường
hợp đơn giản hoặc đặc biệt là một lựa chọn tốt. Các bài xác suất thường có
nhiều trạng thái và khá trừu tượng; chỉ khi quan sát quy luật qua ví dụ cụ thể
ta mới nắm được mô hình và có cơ sở để giải bài toán.

\section{Bài toán về biến ngẫu nhiên liên tục}

Bài toán về biến ngẫu nhiên liên tục chưa xuất hiện nhiều trong thi Tin học,
có lẽ vì lời giải thường phải dùng tích phân. Tuy nhiên kỹ thuật tích phân
trong các bài này thường không phải phần khó; nhiều nhất cũng chỉ dùng công
thức Newton--Leibniz. Vì vậy trong tương lai có khả năng sẽ xuất hiện nhiều
bài hơn về biến ngẫu nhiên liên tục. Các tư tưởng khi giải loại bài này cũng
hữu ích cho những dạng bài khác. Sau đây ta xét hai ví dụ để bàn về các hướng
tiếp cận thường gặp.

\subsection{Ví dụ 3: RNG}

\paragraph{Tóm tắt đề bài.}
Có \(N\) bộ sinh số ngẫu nhiên; bộ thứ \(i\) trả về đều một số thực trong
\([0,R_i]\). Hỏi xác suất để tổng mọi số sinh ra không vượt quá \(b\) là bao
nhiêu? Đáp án cần chính xác tới \(10^{-9}\).

Ràng buộc:
\[
  1\le N\le 10,\qquad
  1\le R_i\le 10,\qquad
  0\le b\le 100,
\]
trong đó \(R_i\) và \(b\) đều là số nguyên.

\paragraph{Phân tích.}
Gọi giá trị trả về của bộ sinh thứ \(i\) là biến ngẫu nhiên \(X_i\). Khi đó
\(X_i\) phân bố đều trên \([0,R_i]\), và \(N\) biến ngẫu nhiên này độc lập với
nhau. Tổng của chúng cũng là một biến ngẫu nhiên; đặt
\[
  S=X_1+X_2+\cdots+X_N.
\]
Ta cần tính \(\Prob(S\le b)\).

Ta sẽ giải bài toán bằng hai cách, lần lượt từ góc nhìn hình học và đại số.

\subsubsection{Cách 1}

Nhớ lại phần kiến thức cơ sở: \(N\) biến ngẫu nhiên tương ứng với không gian
\(N\) chiều, khá phức tạp. Vì vậy trước hết xét các trường hợp \(N\) nhỏ.

Khi \(N=1\), miền giá trị \([0,R_1]\) của \(X_1\) là một đoạn thẳng trên trục
số, còn \(S=X_1\le x\) là một nửa đường thẳng. Xác suất \(\Prob(S\le x)\) là
tỉ số giữa độ dài giao của chúng và \(R_1\).

\begin{center}
\fbox{\parbox{0.82\linewidth}{
Mô tả hình: một trục số có các mốc \(0\), \(x\), \(R_1\); phần từ \(0\) tới
\(x\) biểu diễn miền giao khi \(0\le x\le R_1\).
}}
\end{center}

Khi \(N=2\), miền giá trị của \((X_1,X_2)\) là một hình chữ nhật trong mặt
phẳng tọa độ; điều kiện \(S=X_1+X_2\le x\) là một nửa mặt phẳng. Xác suất
\(\Prob(S\le x)\) là tỉ số giữa diện tích giao của chúng và diện tích hình chữ
nhật \(R_1R_2\).

Khi \(N=3\), miền giá trị của \((X_1,X_2,X_3)\) là một hình hộp chữ nhật trong
không gian; điều kiện \(S=X_1+X_2+X_3\le x\) là một nửa không gian. Xác suất
\(\Prob(S\le x)\) là tỉ số giữa thể tích giao của chúng và thể tích hình hộp
\(R_1R_2R_3\).

\begin{center}
\fbox{\parbox{0.82\linewidth}{
Mô tả hình: một hình hộp chữ nhật trong hệ tọa độ ba chiều bị cắt bởi mặt
phẳng \(x_1+x_2+x_3=x\); phần nằm về phía gốc tọa độ là miền cần lấy thể tích.
}}
\end{center}

Tổng quát, với mọi \(N\), miền giá trị của
\((X_1,X_2,\ldots,X_N)\) là một vùng trong không gian \(N\) chiều, còn
\(S\le x\) là một nửa không gian \(N\) chiều. Tỉ số giữa thể tích \(N\) chiều
của giao của chúng và
\[
  R_1R_2\cdots R_N
\]
chính là \(\Prob(S\le x)\). Vì \(R_1R_2\cdots R_N\) là hằng số, bài toán
chuyển thành tính thể tích của giao. Ký hiệu
\[
  \Vol([a_1,b_1],[a_2,b_2],\ldots,[a_N,b_N],x)
\]
là thể tích \(N\) chiều của miền các điểm có \(X_i\in[a_i,b_i]\) và tổng các
tọa độ không vượt quá \(x\).

Mỗi biến đều có miền giá trị hữu hạn. Khi \(N\) lớn, vùng hình học có rất
nhiều mặt biên và tình huống trở nên phức tạp. Lúc này ta dùng một kỹ thuật
quen thuộc: chuyển sang phần bù,
\[
  [0,R_i]=[0,+\infty)-(R_i,+\infty).
\]
Do đó
\[
  \Vol([0,R_1],[0,R_2],\ldots,[0,R_N],x)
\]
có thể được biểu diễn bằng các thể tích trong đó mỗi \(X_i\) có miền giá trị
\([0,+\infty)\) hoặc \((R_i,+\infty)\). Chi tiết biểu diễn bằng nguyên lý bù
trừ được đặt trong phụ lục.

Bây giờ cần tính thể tích của miền tổng không vượt quá \(x\) khi các biến có
miền giá trị \([0,+\infty)\) hoặc \((R_i,+\infty)\). Với miền
\((R_i,+\infty)\), thay \(X_i\) bằng \(X'_i=X_i-R_i\) và đồng thời giảm \(x\)
đi \(R_i\), bài toán trở thành tương đương, còn miền giá trị của \(X'_i\) là
\([0,+\infty)\).

Vì vậy bài toán quy về tính
\[
  \Vol([0,+\infty),[0,+\infty),\ldots,[0,+\infty),x).
\]
Qua quan sát và chứng minh đơn giản, giá trị này bằng
\[
  \frac{x^N}{N!}
\]
khi \(x\ge 0\), và bằng \(0\) khi \(x<0\). Chứng minh được đưa trong phụ lục.
Kết hợp với bù trừ, ta có công thức trực tiếp:
\[
  \Prob(S\le b)
  =
  \frac{1}{R_1R_2\cdots R_N}
  \sum_{Q\subseteq\{1,2,\ldots,N\}}
  (-1)^{|Q|}
  \frac{\max\left(b-\sum_{i\in Q}R_i,0\right)^N}{N!}.
\]

Đến đây bài toán đã được giải quyết. Nhìn lại quá trình, ta bắt đầu từ các
trường hợp \(N=1,2,3\), đưa bài toán về tính thể tích giao, rồi dùng phần bù
và phép tịnh tiến để quy về thể tích của một đơn hình trong góc phần tư dương.
Qua từng lần biến đổi, hiểu biết về bài toán trở nên sâu hơn.

\subsubsection{Cách 2}

Ta vẫn bắt đầu bằng các trường hợp \(N\) nhỏ.

Khi \(N=1\), xét \(x\) thay đổi trên \([0,+\infty)\), ta có
\[
  \Prob(S\le x)=
  \begin{cases}
    x/R_1, & 0\le x\le R_1,\\
    1, & x>R_1.
  \end{cases}
\]

Khi \(N=2\), \(\Prob(S\le x)\) là diện tích phần của hình chữ nhật nằm dưới
đường thẳng \(x_1+x_2=x\), chia cho diện tích hình chữ nhật.

\begin{center}
\fbox{\parbox{0.82\linewidth}{
Mô tả hình: hình chữ nhật \([0,R_1]\times[0,R_2]\) bị cắt bởi đường thẳng
\(x_1+x_2=x\); phần phía dưới đường thẳng được tô bóng.
}}
\end{center}

Để tính phần nằm dưới đường thẳng, dùng các đường \(x+y=R_1\) và \(x+y=R_2\)
chia hình chữ nhật thành ba phần. Trong mỗi phần, độ dài đoạn mà đường thẳng
\(x_1+x_2=x\) cắt hình chữ nhật biến thiên tuyến tính theo \(x\). Giả sử
\(R_1\ge R_2\), ta thu được
\[
  \Prob(S\le x)=
  \begin{cases}
    \dfrac{x^2}{2R_1R_2}, & 0\le x\le R_2,\\[0.8em]
    \dfrac{2x-R_2}{2R_1}, & R_2<x\le R_1,\\[0.8em]
    \dfrac{-x^2+(2R_1+2R_2)x-(R_1^2+R_2^2)}{2R_1R_2},
      & R_1<x\le R_1+R_2,\\[0.8em]
    1, & R_1+R_2<x.
  \end{cases}
\]

Nhìn lại quá trình tính \(\Prob(S\le x)\) cho \(N=1\) và \(N=2\), ta thấy có
thể chia trục \(x\) thành một số khoảng; trên mỗi khoảng, \(\Prob(S\le x)\) là
một đa thức theo \(x\). Chỉ cần thế \(x\) vào đa thức của khoảng chứa nó là có
kết quả. Vấn đề là phải chia khoảng thế nào.

Vì \(R_1,R_2,\ldots,R_N\) đều là số nguyên, ta tự nhiên đưa ra phỏng đoán:
\(\Prob(S\le x)\) có thể được biểu diễn bằng đa thức trên từng khoảng nguyên
liên tiếp. Phỏng đoán này là đúng. Ta chứng minh như sau.

Gọi \(f_i(x)\) là hàm mật độ xác suất của tổng
\[
  S_i=X_1+X_2+\cdots+X_i.
\]
Trước hết chứng minh rằng \(f_i(x)\) có thể được biểu diễn bằng đa thức trên
từng khoảng nguyên liên tiếp.

Dùng quy nạp. Khi \(i=1\), hàm cần xét chính là mật độ của \(X_1\). Vì \(X_1\)
phân bố đều trên \([0,R_1]\), theo định nghĩa mật độ xác suất,
\[
  f_1(x)=
  \begin{cases}
    1/R_1, & 0\le x\le R_1,\\
    0, & x>R_1,
  \end{cases}
\]
rõ ràng đều là đa thức.

Giả sử kết luận đúng với các \(i=2,3,\ldots,N-1\). Xét \(i=N\). Vì
\(X_1,X_2,\ldots,X_i\) độc lập, tổng \(S_{i-1}\) độc lập với \(X_i\), còn mật
độ của \(X_i\) là hàm từng đoạn nhận giá trị \(1/R_i\) trên \([0,R_i]\) và
\(0\) ở nơi khác. Theo định nghĩa mật độ đồng thời,
\[
  f_i(x)
  =
  \int_{\max(0,x-R_i)}^{x}
  \frac1{R_i}f_{i-1}(t)\dd t.
\]

Lấy tùy ý một khoảng nguyên \([L-1,L]\), với \(L=1,2,\ldots\). Trước hết xét
trường hợp \(\max(0,x-R_i)=x-R_i\). Ta chia khoảng tích phân thành
\[
  [x-R_i,L-R_i], [L-R_i,L-R_i+1],\ldots,[L-2,L-1],[L-1,x].
\]
Ngoài khoảng đầu và khoảng cuối, mọi khoảng đều có hai đầu mút là hằng số, nên
giá trị tích phân là hằng số. Với khoảng đầu và khoảng cuối, dùng công thức
Newton--Leibniz; do nguyên hàm của đa thức vẫn là đa thức, vế phải là tổng của
hai đa thức và một số hằng số, nên vẫn là đa thức.

Nếu \(\max(0,x-R_i)=0\), thay đổi duy nhất là khoảng đầu trở thành \([0,1]\);
kết luận trên vẫn đúng. Vì \([L-1,L]\) được lấy tùy ý, \(f_i(x)\) với \(i=N\)
cũng có thể được biểu diễn bằng đa thức trên từng khoảng nguyên liên tiếp.

Theo định nghĩa hàm mật độ xác suất,
\[
  \Prob(S\le x)
  =
  \int_0^x f_N(t)\dd t
  =
  \int_0^{\lfloor x\rfloor}f_N(t)\dd t
  +
  \int_{\lfloor x\rfloor}^{x}f_N(t)\dd t.
\]
Hạng thứ nhất ở vế phải là hằng số; hạng thứ hai dùng Newton--Leibniz vẫn cho
đa thức. Vì vậy \(\Prob(S\le x)\) có thể được biểu diễn bằng đa thức trên từng
khoảng nguyên liên tiếp.

Đến đây bài toán đã được giải quyết. Ta cũng xuất phát từ các trường hợp
\(N\) nhỏ, đưa ra phỏng đoán rằng \(\Prob(S\le x)\) là đa thức theo đoạn trên
các khoảng nguyên, rồi chứng minh phỏng đoán đó. Bản thân chứng minh đồng thời
cung cấp thuật toán. Trong cách 2, quan sát và phỏng đoán mạnh dạn là cốt lõi
của lời giải.

\subsubsection{So sánh hai cách}

Trong cách 1, công thức thu được tương đối đơn giản. Ở cách 2, để tính
\(\Prob(S\le x)\) ta phải tích phân lặp nhiều lần, phức tạp hơn. Chứng minh
của cách 1 cũng đơn giản hơn, còn cách 2 dùng tới mật độ xác suất đồng thời
nên trừu tượng hơn.

Nhưng cách 2 có ưu điểm là dễ khái quát. Điều kiện mỗi biến ngẫu nhiên phân bố
đều có thể được nới lỏng trong chứng minh của cách 2: chỉ cần hàm mật độ của
mỗi biến ngẫu nhiên là đa thức. Ngược lại, trong cách 1, lý do cốt lõi để có
thể xem xác suất như tỉ số thể tích \(N\) chiều chính là mọi biến đều phân bố
đều.

Chẳng hạn, nếu cần tính xác suất để tổng bình phương của \(N\) số ngẫu nhiên
không vượt quá \(x\), thì cách 1 không còn tác dụng trực tiếp, còn cách 2 chỉ
cần áp dụng tương tự. Cả hai cách đều bắt đầu từ việc quan sát trường hợp đơn
giản, nhưng cách 2 nắm bản chất toán học của bài toán sâu hơn, vì vậy nó trở
thành một phương pháp tổng quát cho bài toán về biến ngẫu nhiên liên tục.

\subsection{Ví dụ 4: Random Shooting}

\paragraph{Tóm tắt đề bài.}
Một đa giác lồi \(N\) đỉnh nằm trong hình vuông có bốn đỉnh
\[
  C(0,0),\quad D(100,0),\quad E(100,100),\quad F(0,100).
\]
Chọn ngẫu nhiên hai điểm trong hình vuông. Hỏi xác suất để chúng thỏa mãn ít
nhất một trong hai điều kiện sau là bao nhiêu?
\begin{enumerate}
  \item Ít nhất một trong hai điểm nằm trong đa giác lồi.
  \item Đoạn thẳng nối hai điểm cắt đa giác lồi.
\end{enumerate}
Ràng buộc: \(3\le N\le 8\), không có ba đỉnh nào của đa giác lồi thẳng hàng,
và tọa độ các đỉnh là các số nguyên dương từ \(1\) tới \(99\).

\paragraph{Phân tích.}
Trước hết cần đơn giản hóa bài toán. Gọi diện tích đa giác lồi là \(S\), diện
tích hình vuông là
\[
  \Total=100\cdot100=10000.
\]
Xác suất hai điểm đều nằm trong đa giác là
\[
  P_1=\left(\frac{S}{\Total}\right)^2.
\]
Xác suất đúng một điểm nằm trong đa giác là
\[
  P_2=2\left(\frac{S}{\Total}\right)\left(1-\frac{S}{\Total}\right).
\]
Đặt
\[
  P_3=\sum_{i=1}^{N}
  \Prob(\text{đoạn nối hai điểm cắt cạnh thứ } i \text{ của đa giác lồi}).
\]

Xét tập các cặp điểm có đoạn nối cắt cạnh thứ \(i\). Chia chúng thành hai
phần.

Phần 1: một điểm ở ngoài đa giác và một điểm ở trong đa giác. Vì đa giác lồi,
trong trường hợp này đoạn nối chỉ có đúng một giao điểm với biên đa giác
(bỏ qua trường hợp cắt đúng đỉnh; tương tự bên dưới). Vì vậy với các cạnh
khác nhau, các phần 1 không giao nhau. Tổng xác suất của phần 1 trên mọi cạnh
chính là \(P_2\).

Phần 2: hai điểm đều ở ngoài đa giác. Vì đa giác lồi, nếu đoạn nối cắt đa giác
thì nó nhất định cắt đúng hai cạnh. Do đó mỗi cặp điểm thuộc phần 2 được tính
đúng hai lần khi cộng trên các cạnh. Vậy tổng xác suất của phần 2 trên mọi
cạnh là
\[
  2\Prob(\text{hai điểm đều ở ngoài đa giác và đoạn nối cắt đa giác}).
\]
Suy ra
\[
  P_3=
  2\Prob(\text{hai điểm đều ở ngoài đa giác và đoạn nối cắt đa giác})+P_2.
\]
Xác suất cần tìm là
\[
  P_1+P_2+\frac{P_3-P_2}{2}.
\]
Như vậy bài toán chuyển thành tính \(P_3\), tức là tính xác suất để đoạn nối
hai điểm cắt từng cạnh. Độ phức tạp của bài toán giảm đi đáng kể.

Xét một cạnh \(AB\) đang cần tính xác suất cắt, ký hiệu xác suất đó là \(P\).
Ta chọn hướng sao cho \(B_y\ge A_y\); nếu \(A_y>B_y\) thì đổi chỗ \(A\) và
\(B\). Như vậy góc giữa đoạn có hướng \(AB\) và trục \(x\) nằm trong
\([0,\pi]\). Lý do là khi đoạn nối hai điểm cắt \(AB\), một điểm nhất định ở
bên trái \(AB\), điểm còn lại ở bên phải \(AB\). Khi đó ta chỉ cần thảo luận
điểm ở bên trái \(AB\).

\begin{center}
\fbox{\parbox{0.82\linewidth}{
Mô tả hình 1: với một điểm cố định \(V(x_0,y_0)\), kẻ hai tia từ \(V\) qua
\(A\) và \(B\). Vùng trong hình vuông bị hai tia này chắn khuất chính là tập
các điểm \(U\) sao cho đoạn \(UV\) cắt \(AB\).
}}
\end{center}

Với điểm cố định \(V(x_0,y_0)\), đặt
\[
  E=\{(x,y)\mid \text{đoạn nối }U(x,y)\text{ với }V\text{ cắt }AB\}.
\]
Tập \(E\) chính là vùng bị tô bóng trong hình minh họa. Có thể tưởng tượng
\(V\) là một nguồn sáng điểm, còn \(AB\) là một tấm chắn; vùng \(V\) không
chiếu tới chính là \(E\). Nói cách khác, kẻ từ \(V\) hai tia qua \(A\) và
\(B\), lấy các giao điểm của chúng với hình vuông, cùng các đỉnh hình vuông và
\(A,B\), ta thu được vùng cần xét. Gọi diện tích vùng bóng ứng với điểm
\(V(x_0,y_0)\) là \(S_V(x_0,y_0)\). Theo định nghĩa phân bố xác suất của biến
ngẫu nhiên liên tục,
\[
  P=\int_0^{100}\int_0^{100}\frac{S_V(x,y)}{\Total}\dd y\dd x.
  \tag{3}
\]
Nhưng \(V\) có vô hạn vị trí, nên không thể tính riêng \(S_V(x,y)\) cho từng
điểm. Ta cần quan sát sâu hơn.

\begin{center}
\fbox{\parbox{0.82\linewidth}{
Mô tả hình 2: từ bốn đỉnh \(C,D,E,F\) của hình vuông kẻ các tia qua điểm
\(A\). Những tia này cắt biên hình vuông tại \(G,H,I,J\), qua đó chia hình
vuông thành tám vùng; trong cùng một vùng, tia từ điểm bất kỳ qua \(A\) luôn
cắt cùng một cạnh của hình vuông.
}}
\end{center}

Xét điểm \(A\) của đoạn có hướng \(AB\). Kẻ các tia từ \(C,D,E,F\) qua \(A\);
chúng lần lượt cắt hình vuông tại \(G,H,I,J\), chia hình vuông thành tám vùng.
Dễ thấy mọi điểm trong cùng một vùng, khi kẻ tia qua \(A\), đều cắt cùng một
cạnh của hình vuông. Ví dụ các điểm trong vùng \(FAC\) khi kẻ tia qua \(A\)
đều cắt cạnh \(ED\); các điểm trong vùng \(GAJ\) đều cắt cạnh \(FC\); các điểm
trong vùng \(HAE\) đều cắt cạnh \(CD\). Gọi \(CD\) là cạnh dưới, \(DE\) là
cạnh phải, \(EF\) là cạnh trên, \(FC\) là cạnh trái. Khi đó mọi điểm trong
cùng một vùng tương ứng với đúng một trong bốn cạnh này.

\begin{center}
\fbox{\parbox{0.82\linewidth}{
Mô tả hình 3: kéo dài \(AB\) tới khi cắt hình vuông tại \(Q,R\). Chỉ xét phần
bên trái của đường có hướng \(AB\), tức tứ giác \(CQRF\). Các tia từ \(D,E\)
qua \(A,B\) chia tứ giác này thành những vùng mà trong mỗi vùng, các tia qua
\(A\) và qua \(B\) cắt hai cạnh cố định của hình vuông.
}}
\end{center}

Quay lại đoạn \(AB\). Kéo dài \(AB\) tới khi cắt hình vuông tại \(Q,R\). Theo
kết luận trên, ta chỉ cần xét phía bên trái của \(QR\), tức tứ giác \(CQRF\).
Kẻ từ \(D,E\) các tia qua \(A,B\); giao điểm của chúng với hình vuông chia
\(CQRF\) thành một số vùng. Trong cùng một vùng, tia từ điểm đó qua \(A\) và
qua \(B\) cắt các cạnh cố định của hình vuông. Có thể phân loại vùng theo cặp
cạnh mà hai tia qua \(A\) và \(B\) cắt. Vì các vùng đều ở bên trái \(AB\), nên
với mọi điểm \(V\), tia \(VB\) nằm theo chiều ngược kim đồng hồ so với tia
\(VA\). Do đó có nhiều nhất 10 loại:
\[
\begin{gathered}
\text{A trái, B trái};\quad \text{A trái, B dưới};\quad
\text{A trái, B phải};\quad \text{A trái, B trên};\\
\text{A dưới, B dưới};\quad \text{A dưới, B phải};\quad
\text{A dưới, B trên};\\
\text{A phải, B phải};\quad \text{A phải, B trên};\quad
\text{A trên, B trên}.
\end{gathered}
\]

\begin{center}
\fbox{\parbox{0.82\linewidth}{
Mô tả hình 4: trong một vùng cố định, khi \(V\) dịch tới \(V'\), giao điểm
của tia \(VA\) với biên hình vuông dịch từ \(Z\) tới \(Z'\). Phần diện tích
thay đổi ở phía \(A\) là tam giác \(ZAZ'\); phía \(B\) có công thức tương tự.
}}
\end{center}

Xét một vùng cụ thể. Gọi \(Z\) là giao điểm của tia \(VA\) với cạnh tương ứng
của hình vuông. Khi \(V\) di chuyển tới \(V'\), giao điểm mới là \(Z'\). Như
hình minh họa, chênh lệch diện tích sinh ra ở phía \(A\) chính là diện tích
tam giác \(ZAZ'\), có thể mang dấu âm. Phía \(B\) cũng có kết luận tương tự.
Vì \(VB\) nằm theo chiều ngược kim đồng hồ so với \(VA\), hai phía này không
cắt nhau và không ảnh hưởng lẫn nhau.

Điều này cho biết, với mọi điểm \(V(x,y)\) trong cùng một vùng, diện tích bóng
\(S_V(x,y)\) bằng một hằng số trừ đi diện tích của hai tam giác. Gọi hằng số
đó là diện tích cơ sở, ta có
\[
  S_V(x,y)=\text{diện tích cơ sở}-f(x,y)-g(x,y),
\]
trong đó \(f(x,y)\) và \(g(x,y)\) lần lượt biểu diễn phần sai khác ở phía
\(A\) và phía \(B\); chúng chỉ phụ thuộc vào loại vùng.

Xét vùng góc trên trái trong hình 4. Vùng này thuộc loại ``A dưới, B phải'',
và diện tích cơ sở có thể xem là diện tích đa giác \(ACDEB\), một đa giác lõm.
Khi đó
\[
\begin{aligned}
f(x,y)
&=\frac12 A_y\left(A_x-A_y\frac{x-A_x}{y-A_y}\right)
 =\frac12\left(A_xA_y-A_y^2\frac{x-A_x}{y-A_y}\right),\\
g(x,y)
&=\frac12(100-B_x)
  \left((100-B_y)-(100-B_x)\frac{y-B_y}{x-B_x}\right)\\
&=\frac12\left((100-B_x)(100-B_y)
-(100-B_x)^2\frac{y-B_y}{x-B_x}\right).
\end{aligned}
\]

Với mỗi vùng, ta có thể tính diện tích cơ sở, \(f\) và \(g\). Vì \(f\) và
\(g\) thực chất đều là diện tích của một tam giác, nên trên mỗi vùng
\(S_V(x,y)\) có dạng
\[
  S_V(x,y)
  =
  c_0+c_1\frac{x-x_0}{y-y_0}
     +c_2\frac{y-y_1}{x-x_1},
  \tag{4}
\]
trong đó \(c_0,c_1,c_2,x_0,y_0,x_1,y_1\) đều là hằng số. Thế (4) vào (3),
\[
\begin{aligned}
P=\frac1{\Total}
\sum_{\text{vùng }Q}
&\Biggl(
c_0\iint_{(x,y)\in Q}\dd x\dd y
c_1\iint_{(x,y)\in Q}\frac{x-x_0}{y-y_0}\dd x\dd y\\
&\qquad\qquad
c_2\iint_{(x,y)\in Q}\frac{y-y_1}{x-x_1}\dd y\dd x
\Biggr).
\end{aligned}
\tag{5}
\]

Để tính
\[
  \iint_{(x,y)\in Q}\frac{y-y_1}{x-x_1}\dd y\dd x,
\]
ta chia mỗi vùng bằng các đường thẳng đứng thành một số mảnh nhỏ; mỗi mảnh là
một tứ giác có hai cạnh thẳng đứng. Giả sử bốn đỉnh của một tứ giác như vậy là
\[
  (\mathit{st},\mathit{bot}_{st}),\quad
  (\mathit{st},\mathit{top}_{st}),\quad
  (\mathit{en},\mathit{bot}_{en}),\quad
  (\mathit{en},\mathit{top}_{en}),
\]
trong đó
\[
  \mathit{st}<\mathit{en},\qquad
  \mathit{bot}_{st}\le \mathit{top}_{st},\qquad
  \mathit{bot}_{en}\le \mathit{top}_{en}.
\]
Hai cạnh dưới và trên lần lượt có thể viết là
\[
  y=a_1x+r_1,\qquad y=a_2x+r_2,
\]
với \(a_1,a_2,r_1,r_2\) là hằng số. Khi đó
\[
\begin{aligned}
\iint_Q\frac{y-y_1}{x-x_1}\dd y\dd x
&=
\int_{\mathit{st}}^{\mathit{en}}
\int_{a_1x+r_1}^{a_2x+r_2}
\frac{y-y_1}{x-x_1}\dd y\dd x\\
&=
\int_{\mathit{st}}^{\mathit{en}}
\frac{ax^2+bx+c}{x-x_1}\dd x\\
&=
\int_{\mathit{st}}^{\mathit{en}}
\left(a'x+b'+\frac{c'}{x-x_1}\right)\dd x\\
&=
\frac{a'(\mathit{en}^2-\mathit{st}^2)}{2}
+b'(\mathit{en}-\mathit{st})\\
&\quad
+c'\bigl(\ln|\mathit{en}-x_1|-\ln|\mathit{st}-x_1|\bigr),
\end{aligned}
\]
trong đó \(a,b,c,a',b',c'\) đều là hằng số. Tổng giá trị tích phân trên các tứ
giác nhỏ chính là tích phân trên cả vùng \(Q\).

Để tính
\[
  \iint_{(x,y)\in Q}\frac{x-x_0}{y-y_0}\dd x\dd y,
\]
chỉ cần chia vùng bằng các đường ngang thành các tứ giác có hai cạnh ngang rồi
tính bằng phương pháp tương tự.

Đến đây bài toán đã được giải quyết. Nhìn toàn bộ quá trình, chính nhờ liên
tục biến đổi mà độ phức tạp của bài toán giảm xuống mức có thể xử lý. Quan
sát và phân tích tỉ mỉ là chìa khóa; kỹ thuật tích phân được dùng không phức
tạp và không phải điểm khó. Điều này gợi ý rằng khi giải bài toán phân bố xác
suất của biến ngẫu nhiên liên tục, nền tảng toán học là cơ sở, còn quan sát và
phân tích mới là then chốt.

\section{Tổng kết}

Bài viết đã xét hai ví dụ liên quan tới biến ngẫu nhiên rời rạc và hai ví dụ
liên quan tới biến ngẫu nhiên liên tục. Ở ví dụ thứ nhất, ta linh hoạt dùng
tính chất của tổ hợp để giải bài toán xác suất thắng tối đa. Ở ví dụ thứ hai,
ta định nghĩa các tầng của bộ sinh số ngẫu nhiên và dùng phân tích sai số để
thu được thiết bị kích hoạt ngẫu nhiên tối ưu. Ở ví dụ thứ ba, hai cách dùng
chuyển sang phần bù và chia hàm theo đoạn lần lượt giải bài toán phân bố của
tổng nhiều biến ngẫu nhiên, rồi được đem ra so sánh. Ở ví dụ thứ tư, thông
qua đơn giản hóa đa giác lồi và chia trường hợp, cuối cùng ta đưa bài toán về
tính một biểu thức tích phân.

Bốn lời giải có đặc sắc riêng, phản ánh tính phức tạp và biến hóa của bài
toán xác suất. Tuy vậy, từ quá trình phân tích cũng có thể rút ra một vài
nguyên tắc chung. Cả bốn bài đều yêu cầu nắm chính xác các định nghĩa toán học
và kiến thức liên quan; đều bắt đầu bằng phân tích các trường hợp đơn giản và
nhấn mạnh việc giản lược bài toán. Các nguyên tắc này xuất phát từ đặc điểm
của bài toán xác suất: tính toán học mạnh, trừu tượng và phức tạp. Biến phức
tạp thành đơn giản, biến trừu tượng thành trực quan, là tư tưởng quan trọng
khi phân tích và giải bài toán xác suất.

\section*{Lời cảm ơn}

Xin cảm ơn thầy Liu Rujia đã hướng dẫn lựa chọn đề tài.
Xin cảm ơn bạn Shou Heming đã đưa ra các ý kiến quý báu khi sửa bài viết.
Xin cảm ơn bạn Jin Bin của đội tuyển tập huấn đã hỗ trợ trong việc giải bài.
Xin cảm ơn bạn Tang Keyin của đội tuyển tập huấn đã giúp đỡ.
Xin cảm ơn các huấn luyện viên Tang Wenbin và Hu Weidong đã chỉ đạo.

\section*{Tài liệu tham khảo}

\begin{enumerate}
  \item Thomas H. Cormen và cộng sự, \emph{Introduction to Algorithms}.
  \item He Shuyuan, \emph{Xác suất luận}.
  \item Liu Rujia và Huang Liang, \emph{Nghệ thuật thuật toán và thi Tin học}.
  \item Chang Gengzhe và Shi Jihuai, \emph{Giáo trình giải tích toán học}.
  \item Các bài viết của đội tuyển tập huấn quốc gia từ năm 2005 tới 2008.
\end{enumerate}

\appendix

\section{Biểu diễn thể tích vùng}

Ta cần biểu diễn
\[
  \Vol([0,R_1],[0,R_2],\ldots,[0,R_N],x).
\]
Dùng
\[
  [0,R_i]=[0,+\infty)-(R_i,+\infty),
\]
ta có thể viết hình thức
\[
\begin{aligned}
&\Vol([0,R_1],[0,R_2],\ldots,[0,R_N],x)\\
&=
\Vol([0,+\infty)-(R_1,+\infty),
     [0,+\infty)-(R_2,+\infty),
     \ldots,
     [0,+\infty)-(R_N,+\infty),x).
\end{aligned}
\]
Đặt \(D_i\in\{0,R_i\}\) với \(1\le i\le N\). Gọi
\[
  Q=\{i\mid D_i=R_i\}.
\]
Khi đó
\[
\begin{aligned}
\Vol([0,R_1],[0,R_2],\ldots,[0,R_N],x)
&=
\sum_{Q\subseteq\{1,2,\ldots,N\}}
(-1)^{|Q|}\\
&\quad\cdot
\Vol((D_1,+\infty),(D_2,+\infty),\ldots,(D_N,+\infty),x).
\end{aligned}
\]
Công thức này gần giống công thức của nguyên lý bù trừ.

\section{Chứng minh công thức thể tích trong cách 1 của ví dụ RNG}

Ta chứng minh
\[
  \Vol([0,+\infty),[0,+\infty),\ldots,[0,+\infty),x)
  =\frac{x^N}{N!}
\]
với \(x\ge 0\).

Dùng quy nạp theo \(N\). Khi \(N=1\), ta có \(\Vol([0,+\infty),x)=x\), hiển
nhiên đúng.

Giả sử kết luận đúng với \(N=2,3,\ldots,k-1\). Khi \(N=k\),
\[
  \Vol([0,+\infty),[0,+\infty),\ldots,[0,+\infty),x)
\]
là thể tích \(k\) chiều của một khối nón \(k\) chiều. Diện tích đáy là
\[
  \frac{x^{k-1}}{(k-1)!}.
\]
Mặt khác, thể tích bằng tích phân của diện tích thiết diện. Với thiết diện
cách đỉnh khối nón một khoảng \(h\), diện tích thiết diện là
\[
  \frac{x^{k-1}}{(k-1)!}\left(\frac{h}{x}\right)^{k-1}.
\]
Do đó
\[
\begin{aligned}
\Vol([0,+\infty),[0,+\infty),\ldots,[0,+\infty),x)
&=
\int_0^x
\frac{x^{k-1}}{(k-1)!}\left(\frac{h}{x}\right)^{k-1}\dd h\\
&=
\frac{1}{(k-1)!}\left(\frac{x^k}{k}-0\right)
=\frac{x^k}{k!}.
\end{aligned}
\]
Quy nạp hoàn tất.

\section{Nguyên đề các bài toán}

\subsection*{LastMarble}

\paragraph{Mô tả bài toán.}
Một túi chứa một số viên bi đỏ và xanh. Mỗi lượt, người chơi lấy khỏi túi
\(1\), \(2\) hoặc \(3\) viên bi. Người chơi nhìn màu các viên đã lấy và công
bố số lượng từng màu. Người lấy viên bi đỏ cuối cùng sẽ thua. Cho hai số
nguyên \(\mathit{red}\) và \(\mathit{blue}\), là số bi đỏ và xanh ban đầu.
Trước khi chơi, một người bạn lấy ngẫu nhiên một số viên khỏi túi mà không cho
ai xem; số viên này là \(\mathit{removed}\). Giả sử bạn đi trước và hai bên
đều chọn tối ưu số bi cần lấy, hãy tính xác suất bạn thắng.

\paragraph{Định nghĩa giao diện.}
\[
\begin{array}{ll}
\text{Class:} & \texttt{LastMarble}\\
\text{Method:} & \texttt{winningChance}\\
\text{Parameters:} & \texttt{int, int, int}\\
\text{Returns:} & \texttt{double}\\
\text{Signature:} & \texttt{double winningChance(int red, int blue, int removed)}
\end{array}
\]
Phương thức phải là \texttt{public}.

\paragraph{Ghi chú và ràng buộc.}
Giá trị trả về phải có sai số tuyệt đối hoặc tương đối không quá \(10^{-9}\).
Các ràng buộc là
\[
  1\le \mathit{red}\le 100,\qquad
  1\le \mathit{blue}\le 100,\qquad
  0\le \mathit{removed}\le \mathit{red}-1.
\]

\subsection*{Randomness}

\paragraph{Mô tả bài toán.}
Sinh số ngẫu nhiên thực sự một cách hiệu quả là một vấn đề thú vị trong khoa
học máy tính. Các bộ sinh số ngẫu nhiên khá nhanh trong thư viện thường chỉ là
giả ngẫu nhiên: chúng có chu kỳ hữu hạn và các bit thấp có thể dự đoán được ở
một mức nào đó.

Ngay cả khi bản thân bộ sinh số là hoàn hảo, cách dùng phổ biến của chúng vẫn
có thể không hoàn hảo. Ví dụ, nếu có một hàm trả về đều một số nguyên thật sự
ngẫu nhiên từ \(1\) tới \(100\), và muốn chọn một trong ba sự kiện với xác
suất \(1/3\) cho mỗi sự kiện, ta không thể làm điều đó chỉ bằng một lần gọi.

Tuy nhiên, nếu cho phép gọi thêm bộ sinh khi cần, ta có thể đạt phân bố đều
hoàn hảo trên ba sự kiện. Thực tế, bất kỳ bộ sinh hoàn hảo nào cũng có thể
được dùng để chọn trong bất kỳ tập sự kiện nào với các xác suất cho trước một
cách chính xác.

Giả sử ta có một bộ sinh hoàn hảo trả về số từ \(1\) tới \(R\), và có các xác
suất mong muốn cho \(N\) sự kiện rời nhau. Hãy xác định một thuật toán dùng bộ
sinh này để chọn một sự kiện với xác suất chính xác, đồng thời tối thiểu hóa
số lần gọi kỳ vọng.

\paragraph{Dữ liệu vào.}
Dữ liệu gồm nhiều nhất \(32\) bộ test, mỗi bộ gồm hai dòng. Dòng đầu chứa hai
số nguyên \(R,N\), với
\[
  2\le R\le 1000,\qquad 2\le N\le 1000.
\]
Dòng thứ hai chứa \(N\) cặp số nguyên \(a_i,b_i\), với
\[
  1\le a_i<b_i\le 1000.
\]
Xác suất mong muốn của sự kiện \(i\) là \(a_i/b_i\), và tổng các xác suất bằng
\(1\). Dữ liệu kết thúc bằng một dòng gồm hai số \(0\).

\paragraph{Dữ liệu ra.}
Với mỗi bộ test, in ra số lần gọi kỳ vọng nhỏ nhất cần thiết, làm tròn tới sáu
chữ số sau dấu thập phân.

\subsection*{RNG}

\paragraph{Mô tả bài toán.}
LoadingTime nhận được một bộ sinh số ngẫu nhiên từ bạn học và dành nhiều thời
gian nghiên cứu nó. Bộ sinh này có thể tạo một số thực trong đoạn \([-S,S]\)
như sau. Trước hết nó sinh \(n\) số nguyên \(X_1,\ldots,X_n\) có tổng bằng
\(S\). Sau đó, với mỗi \(X_i\), nó sinh ngẫu nhiên đều một số thực trong đoạn
\([-X_i,X_i]\). Kết quả cuối cùng là tổng của \(N\) số thực đã sinh. LoadingTime
nhận thấy phân bố của kết quả rất thú vị và muốn biết: với \(N\) và \(X\) cho
trước, xác suất để số sinh ra nằm trong đoạn \([A,B]\) là bao nhiêu?

\paragraph{Dữ liệu vào.}
Dòng đầu chứa số nguyên \(T\), là số test. Với mỗi test, dòng đầu chứa ba số
nguyên
\[
  N,A,B,\qquad 1\le N\le 10,\quad -100\le A\le B\le 100.
\]
Dòng thứ hai chứa \(X_1,\ldots,X_n\), với \(1\le X_i\le 10\).

\paragraph{Dữ liệu ra.}
Với mỗi test, in một dòng chứa xác suất cần tính, đúng chín chữ số sau dấu
thập phân.

\paragraph{Ví dụ.}
\[
\begin{array}{rcl}
\text{input} &&
\begin{array}{l}
5\\
1\ -100\ 100\\
10\\
1\ 10\ 90\\
10\\
1\ -20\ 5\\
10\\
2\ -20\ 5\\
5\ 5\\
5\ -5\ 10\\
1\ 2\ 3\ 4\ 5
\end{array}
\\[1em]
\text{output} &&
\begin{array}{l}
1.000000000\\
0.000000000\\
0.750000000\\
0.875000000\\
0.864720052
\end{array}
\end{array}
\]

\subsection*{Random Shooting}

\paragraph{Mô tả bài toán.}
Bạn đang chơi một trò mới ở trường bắn. Trước hết, một mục tiêu được đặt trên
một bảng hình vuông. Bản thân mục tiêu là một đa giác lồi nằm hoàn toàn bên
trong bảng.

Sau đó bạn được bắn hai phát. Nếu ít nhất một phát bắn rơi vào trong mục tiêu
thì bạn thắng. Ngoài ra, nếu đoạn nối hai điểm bắn cắt mục tiêu thì bạn vẫn
thắng. Nếu không thỏa điều kiện nào, bạn thua.

Giả sử khả năng ngắm của bạn rất tệ, tức là hai điểm bắn độc lập và phân bố
đều trên toàn bộ bảng. Hãy tính xác suất thắng.

\paragraph{Dữ liệu vào.}
Dòng đầu chứa số nguyên \(N\), với \(3\le N\le 8\), là số đỉnh của mục tiêu.
\(N\) dòng tiếp theo, mỗi dòng chứa hai số nguyên \(x_i,y_i\), với
\[
  1\le x_i,y_i\le 99,
\]
là tọa độ các đỉnh của mục tiêu theo thứ tự ngược chiều kim đồng hồ. Không có
ba đỉnh nào của mục tiêu thẳng hàng. Hệ tọa độ được chọn sao cho bốn góc của
bảng là \((0,0)\), \((0,100)\), \((100,0)\) và \((100,100)\).

\paragraph{Dữ liệu ra.}
In xác suất cần tính. Lời giải được chấp nhận nếu sai số không quá \(10^{-7}\).

\paragraph{Ví dụ.}
\[
\begin{array}{rl}
\text{sample input} &
\begin{array}{l}
4\\
25\ 25\\
75\ 25\\
75\ 75\\
25\ 75
\end{array}
\\[1em]
\text{sample output} &
\begin{array}{l}
0.7328341830
\end{array}
\end{array}
\]

\end{document}
